\subsection{特殊点}

设 L 是属于 W 的平移所成的集合，并设 \(\Lambda\) 为由这样的 \(t \in T\) 组成的集合：平移 \(x \mapsto t + x\) 属于 L。立即可见 \(\Lambda\) 在 U(W) 下稳定，且 L 是 W 的正规子群。由于 W 适当地作用于 E，L 亦然，从而容易得出 \(\Lambda\) 是 T 的离散子群。对 E 的任意一点 x，记 \(W_x\) 为 x 在 W 中的稳定子。

命题 9. 设 \(a \in E\) 。下列条件等价：

\begin{enumerate}
\def\labelenumi{(\roman{enumi})}
\item
  我们有 \(\mathrm{W} = \mathrm{W}_{a}.\mathrm{L}\) ；
\item
  同态 U 在 \(W_{a}\) 上的限制是从 \(W_{a}\) 到 \(\mathrm{U}(\mathrm{W})\) 的同构；
\item
  对每个超平面 \(\mathbf{H} \in \mathfrak{H}\) ，存在超平面 \(\mathbf{H}' \in \mathfrak{H}\) 平行于 \(\mathbf{H}\) 且使得 \(a \in \mathbf{H}'\) 。
\end{enumerate}

显然 (i) \(\Leftrightarrow\) (ii)，因为 L 是 U 的核且 \(\mathrm{L} \cap \mathrm{W}_a = \{1\}\) 。

假设 (i)，并设 \(\mathbf{H} \in \mathfrak{H}\) ；则 \(s_{\mathbf{H}} \in \mathbf{W}_a\) . L，于是存在向量 \(t \in \Lambda\) 使得 \(a = s_{\mathbf{H}}(a) + t\) ；向量 \(t\) 与 \(\mathbf{H}\) 正交，且若 \(\mathrm{H}' = \mathrm{H} + \frac{1}{2} t\) ，则 \(s_{\mathbf{H}'}(x) = s_{\mathbf{H}}(x) + t\) （对所有 \(x \in \mathbf{E}\) ；参见~§2，第 4 号，命题 5）。由于 \(t \in \Lambda\) 且 \(s_{\mathbf{H}} \in \mathbf{W}\) ，我们有 \(s_{\mathbf{H}'} \in \mathbf{W}\) ，从而 \(\mathrm{H}' \in \mathfrak{H}\) ；我们还有 \(a = s_{\mathbf{H}'}(a)\) ，从而 \(a \in \mathbf{H}'\) 。于是 (i) 蕴含 (iii)。

假设 (iii)。设 \(\mathrm{H} \in \mathfrak{H}\) ；取 (iii) 中的 \(\mathrm{H}'\) 。则 \(s_{\mathrm{H}'}(a) = a\) ，故 \(s_{\mathrm{H}'} \in \mathrm{W}_a\) ；由于 \(\mathrm{H}\) 平行于 \(\mathrm{H}'\) ，元素 \(w = s_{\mathrm{H}'} s_{\mathrm{H}}\) （\(\mathrm{W}\) 中的）是一个平移（§2，第 4 号，命题 5），故 \(w \in \mathrm{L}\) ；于是 \(s_{\mathrm{H}} = s_{\mathrm{H}'} w \in \mathrm{W}_a\) . L。由于 \(\mathrm{W}\) 由族 \((s_{\mathrm{H}})_{\mathrm{H} \in \mathfrak{H}}\) 生成，可得 \(\mathrm{W} = \mathrm{W}_a\) . L，因而 (iii) 蕴含 (i)。

定义 1. 一点 \(a\) （属于 \(E\) ）若满足命题 9 中的等价条件，就称为 \(W\) 的特殊点。

显然 E 的特殊点所成的集合在 W 下稳定。

命题 10. 存在 W 的一个特殊点。

由第 8 号命题 6，只需考虑 W 是本质的情形。T 的自同构群 U(W) 是有限的（参见~第 6 号，定理 3），且对每个超平面 H，\(U(s_{H})\) 是一个正交反射；此外，U(W) 由族 \((U(s_{H}))_{H \in \mathfrak{H}}\) 生成。由第 9 号命题 7，存在 T 的一组基 \((e_{i})_{i \in I}\) ，使得群 U(W) 由按下式定义的反射集 \((s_{i})_{i \in I}\) 生成：

\[
s _ {i} (t) = t - 2 (t | e _ {i}). e _ {i}.
\]

第 2 号定理 1 的推论表明，每个反射 \(s \in \mathrm{U}(\mathrm{W})\) 都形如 \(s = \mathrm{U}(s_{\mathrm{H}})\) ，其中 \(\mathrm{H} \in \mathfrak{H}\) 。于是可以在 \(\mathfrak{H}\) 中找到一族超平面 \((\mathbf{H}_i)_{i \in I}\) ，使得 \(s_i = \mathrm{U}(s_{\mathrm{H}_i})\) 对所有 \(i\) 成立。由于诸向量 \(e_i\) 线性无关，存在 \(a \in E\) 使得 \(a \in H_i\) 对所有 \(i \in I\) 成立。我们有 \(s_{\mathrm{H}_i} \in W_a\) ，故 \(\mathrm{U}(\mathrm{W}) = \mathrm{U}(\mathrm{W}_a)\) ，这意味着 \(W = W_a.L\) ，因为 \(L\) 是 \(U\) 的核。于是 \(a\) 是特殊点。

当 W 有限且本质时，W 只有唯一的特殊点，即在 W 下不变的唯一点。因此，对特殊点的考察主要在 W 无限时才有意义。

命题 11. 假设 W 是本质的。设 a 是 W 的一个特殊点。相对于 \(W_{a}\) 的诸室是以 a 为顶点的开单形锥。对每个室 \(C'\) （相对于 \(W_{a}\) ），存在唯一一个相对于 W 的室 C，它包含于 \(C'\) 且满足 \(a \in \overline{C}\) 。诸 \(w'(\overline{C})\) （\(w' \in W_{a}\) ）的并是 a 在 E 中的一个闭邻域。\(C'\) 的每面墙都是 C 的墙。若 W 无限且不可约，则 C 的墙是 \(C'\) 的墙加上一个不与 \(C'\) 的墙平行的仿射超平面。

设 \(\mathfrak{H}'\) 为满足 \(H \in \mathfrak{H}\) 且包含 \(a\) 的 H 所成的集合。群 \(W_a\) 由诸 \(s_H\) （\(H \in \mathfrak{H}'\) ）生成（第 3 号，命题 2）。相对于 \(W_a\) 的诸室是以 \(a\) 为顶点的开单形锥（第 9 号，命题 7）。设 \(C'\) 是这样一个室，\(U\) 是以 \(a\) 为中心、不与 \(\mathfrak{H} - \mathfrak{H}'\) 中任何元素相交的非空开球。由于 \(a \in \overline{C'}\) ，存在一点 \(b\) 属于 \(U \cap C'\) 。于是 \(b \notin H\) 对所有 \(H \in \mathfrak{H}\) 成立，从而 \(b\) 属于一个室 \(C\) （相对于 \(\mathfrak{H}\) ）。由于 \(\mathfrak{H}' \subset \mathfrak{H}\) ，我们有 \(C \subset C'\) 。集合 \(U \cap C'\) 不与任何 \(H \in \mathfrak{H}\) 相交且是凸的，故 \(U \cap C' \subset C\) ；于是 \(a \in \overline{C}\) 。反之，设 \(C_1\) 是相对于 \(W\) 、包含于 \(C'\) 且满足 \(a \in \overline{C}_{1}\) 的一个室；则 \(C_{1}\) 与 U 相交，且 \(U \cap C_{1} \subset U \cap C' = U \cap C\) ；室 C 与 \(C_{1}\) 有公共点，故必重合。对任意 \(w' \in W_{a}\) ，我们有 \(w'(U) = U\) ，于是

\[
\mathrm{U} \cap w ^ {\prime} (\mathrm{C}) = w ^ {\prime} (\mathrm{U} \cap \mathrm{C}) = w ^ {\prime} (\mathrm{U} \cap \mathrm{C} ^ {\prime}) = \mathrm{U} \cap w ^ {\prime} (\mathrm{C} ^ {\prime});
\]

由于诸 \(w'(C)\) （\(w' \in W_a\) ）的并在 \(E\) 中稠密，诸 \(U \cap w'(C) = U \cap w'(C')\) 的并在 \(U\) 中稠密，从而诸 \(w'(\overline{C})\) （\(w' \in W_a\) ）的并包含 \(U\) 。最后，若 \(H\) 是 \(C'\) 的一面墙，则存在一点 \(c \in U \cap H\) 以及开邻域 \(V \subset U\) （\(c\) 的），使得 \(V \cap C'\) 是 \(V\) 与由 \(H\) 界定且包含 \(C'\) 的开半空间之交；由于 \(V \cap C' = V \cap U \cap C' = V \cap U \cap C = V \cap C\) ，可知 \(H\) 是 \(C\) 的一面墙。若 \(W\) 无限且不可约，则 \(C\) 是一个开单形（第 9 号，命题 8），因而比开单形锥 \(C'\) 多一面墙。

推论. 假设 W 是本质的。

\begin{enumerate}
\def\labelenumi{(\roman{enumi})}
\item
  若 \(a \in \mathbb{E}\) 是特殊的，则存在一个室 \(\mathbb{C}\) 使得 \(a\) 是 \(\overline{\mathbb{C}}\) 的一个极点。
\item
  若 \(\mathrm{C}\) 是一个室，则存在 \(\overline{\mathrm{C}}\) 的一个极点，它是特殊点。
\end{enumerate}

第一个断言由命题 11 得出。第二个断言由第一个以及 W 传递地作用于室集这一事实得出。

然而，\(\overline{C}\) 的极点未必是 W 的特殊点（参见~第 VI 章，图版 X，\(B_{2}\) 型与 \(G_{2}\) 型）。

注. 1) 假设 W 本质、不可约且无限，并沿用命题 11 的记号。由于 U 是从 \(W_{a}\) 到 \(U(W)\) 的同构，可知位移群 \(U(W)\) （它由诸反射 \(U(s_{H})\) （\(H \in \mathfrak{H}\) ）生成）的 Coxeter 图，可由 W 的 Coxeter 图删去顶点 i（它对应于 C 的唯一一面不是 \(C'\) 的墙）而得到。

命题 12. 假设 W 是本质的。设 a 是一个特殊点，L(a) 是 a 在平移群 L 下的诸像所成的集合，并设 C 是一个室。则 \(\overline{C}\) 与 L(a) 交于唯一点；该点是 \(\overline{C}\) 的一个极点。

存在一个室 \(C_{1}\) 使得 a 是 \(\overline{C}_{1}\) 的极点（命题 11 的推论）。每个室都形如 \(\mathbf{C}=tw^{\prime}(\mathbf{C}_{1})\) ，其中 \(w^{\prime}\in W_{a}\) 且 \(t\in L\) ，因为 \(W=W_{a}.L\) ；于是 \(tw^{\prime}(a)=t(a)\in\mathrm{L}(a)\) 是 \(\overline{C}\) 的一个极点。另一方面，\(\overline{C}\) 不可能包含 \(\mathrm{L}(a)\) 中两个不同的点，因为 \(\overline{C}\) 是 W 的基本域（第 3 号，定理 2）。

注. 2) 集合 L(a) 包含于特殊点集，但一般与它不同（参见~第 VI 章，§2，第 2 号以及图版 I 至 VI）。

